\documentclass[12pt]{article}
\usepackage[utf8]{inputenc}
\usepackage{amsmath, amssymb, amsthm}
\usepackage[margin=1in]{geometry}

\newtheorem{theorem}{Theorem}
\newtheorem{lemma}[theorem]{Lemma}

\title{Problem 294: Sum of Digits --- Experience \#23}
\author{Project Euler Solution}
\date{}

\begin{document}
\maketitle

\section{Problem Statement}
For a positive integer $k$, let $d(k)$ be the digit sum of $k$. Define
\[
S(n) = \#\{k : 0 < k < 10^n,\; 23 \mid k,\; d(k) = 23\}.
\]
Given: $S(9) = 263626$, $S(42) = 6\,377\,168\,878\,570\,056$. Find $S(11^{12}) \bmod 10^9$.

\section{Mathematical Foundation}

\begin{theorem}[Digit DP State Space]
The count $S(n)$ equals the $(23, 0)$-entry of the vector $T^n \mathbf{v}_0$, where $T$ is a $552 \times 552$ transition matrix over the state space $\{0,\ldots,23\} \times \{0,\ldots,22\}$ and $\mathbf{v}_0$ has a $1$ at state $(0,0)$.
\end{theorem}

\begin{proof}
Encode partial numbers by $(s, r)$ where $s$ is the running digit sum and $r$ is the running value modulo 23. The state space has $24 \times 23 = 552$ elements. Appending digit $d$ transitions $(s, r) \to (s+d, (10r+d) \bmod 23)$ when $s + d \leq 23$.

The matrix $T$ encodes these transitions: $T_{(s',r'),(s,r)} = \#\{d \in \{0,\ldots,9\} : s' = s+d, r' = (10r+d) \bmod 23\}$. After $n$ digit positions, $S(n) = (T^n)_{(23,0),(0,0)}$. The case $k = 0$ has $d(0) = 0 \neq 23$, so it is automatically excluded.
\end{proof}

\begin{theorem}[Matrix Exponentiation]
For any $n \times n$ matrix $T$ and positive integer $N$, $T^N$ is computable in $O(n^3 \log N)$ arithmetic operations via repeated squaring.
\end{theorem}

\begin{proof}
Write $N = \sum_{i=0}^{\lfloor \log_2 N \rfloor} b_i 2^i$. Compute successive squares $T, T^2, T^4, \ldots$ and multiply those corresponding to $b_i = 1$. This requires $O(\log N)$ matrix multiplications, each costing $O(n^3)$.
\end{proof}

\begin{lemma}[Modular Correctness]
Reducing each matrix entry modulo $10^9$ at every multiplication step yields the correct value of $S(11^{12}) \bmod 10^9$.
\end{lemma}

\begin{proof}
By the homomorphism property of modular arithmetic: $(AB)_{ij} \bmod m = \left(\sum_k (A_{ik} \bmod m)(B_{kj} \bmod m)\right) \bmod m$. Induction on the number of multiplications gives the result.
\end{proof}

\section{Algorithm}

\begin{verbatim}
1. Build 552 x 552 transition matrix T
   State (s, r) indexed as s*23 + r
   For each digit d in {0..9}: T[s+d, (10r+d) mod 23] += 1
2. Compute T^(11^12) mod 10^9 via repeated squaring
3. Return entry at (s=23, r=0) from initial state (s=0, r=0)
\end{verbatim}

\section{Complexity Analysis}
\begin{itemize}
\item \textbf{Time:} $O(552^3 \cdot \log_2(11^{12})) \approx O(552^3 \cdot 42) \approx 7.1 \times 10^9$ modular multiplications.
\item \textbf{Space:} $O(552^2) \approx 3 \times 10^5$ entries for the matrix.
\end{itemize}

\section{Answer}

\[
\boxed{789184709}
\]

\end{document}
